% Bagian: incompleteness
% Bab: theories-computability
% Subbagian: consis-with-q

\documentclass[../../../include/open-logic-section]{subfiles}

\begin{document}

\olfileid[id]{inc}{tcp}{con}

\olsection{Teori yang Konsisten dengan $\Th{Q}$ Tidak Dapat Diputuskan}

Teorema berikut menyatakan bahwa bukan hanya $\Th{Q}$ yang tidak dapat
diputuskan, tetapi sesungguhnya setiap teori yang tidak bertentangan
dengan~$\Th{Q}$ juga tidak dapat diputuskan.

\begin{thm}
Misalkan $\Th{T}$ sebarang teori dalam bahasa aritmetika yang konsisten
dengan~$\Th{Q}$ (yakni $\Th{T} \cup \Th{Q}$ konsisten). Maka $\Th{T}$ tidak
dapat diputuskan.
\end{thm}

\begin{proof}
Ingat bahwa $\Th{Q}$ mempunyai himpunan aksioma berhingga, $!Q_1$, \dots,
$!Q_8$. Kita bahkan dapat mengganti semua aksioma tersebut dengan satu
aksioma, $!E = !Q_1 \land \dots \land !Q_8$.

Andaikan $\Th{T}$ suatu teori yang dapat diputuskan dan konsisten
dengan~$\Th{Q}$. Misalkan
\[
C = \Setabs{!A}{\Th{T} \Proves !E \lif !A}.
\]
Kita tunjukkan bahwa $C$ akan menjadi pemisah dapat dihitung bagi $\Th{Q}$
dan~$\Th{\bar Q}$, suatu kontradiksi. Pertama, jika $!A$ berada
dalam~$\Th{Q}$, maka $!A$ dapat dibuktikan dari aksioma-aksioma~$\Th{Q}$;
menurut teorema deduksi, terdapat !!a{derivation} atas $!E \lif !A$ dalam
logika orde pertama. Jadi $!A$ berada dalam~$C$.

Di pihak lain, jika $!A$ berada dalam~$\Th{\bar Q}$, maka terdapat suatu
pembuktian atas $!E \lif \lnot !A$ dalam logika orde pertama. Jika
$\Th{T}$ juga membuktikan $!E \lif !A$, maka $\Th{T}$ membuktikan
$\lnot !E$; dalam hal itu $\Th{T} \cup \Th{Q}$ tak konsisten. Namun, kita
mengasumsikan bahwa $\Th{T} \cup \Th{Q}$ konsisten, sehingga $\Th{T}$
tidak membuktikan $!E \lif !A$, dan dengan demikian $!A$ tidak berada
dalam~$C$.

Kita telah menunjukkan bahwa jika $!A$ berada dalam~$\Th{Q}$, maka kalimat
itu berada dalam~$C$, dan jika $!A$ berada dalam~$\Th{\bar Q}$, maka kalimat
itu berada dalam~$\Complement{C}$. Jadi $C$ merupakan pemisah dapat dihitung,
dan inilah kontradiksi yang dicari.
\end{proof}

Teorema ini sangat kuat. Sebagai contoh, teorema itu mengimplikasikan:

\begin{cor}
  Logika orde pertama untuk bahasa aritmetika (yakni himpunan
  $\Setabs{!A}{\text{$!A$ dapat dibuktikan dalam logika orde pertama}}$)
  tidak dapat diputuskan.
\end{cor}

\begin{proof}
Logika orde pertama merupakan himpunan konsekuensi dari~$\emptyset$,
yang konsisten dengan~$\Th{Q}$.
\end{proof}

\end{document}
